\documentclass{article}

\usepackage{amsmath}
\usepackage{amsthm}
\usepackage{tikz}
\usepackage{hyperref}

\newtheorem{theorem}{Theorem}

\title{LaTeX Visual Editor}
\author{Visual Editor Example}
\date{\today}

\begin{document}

\maketitle

\begin{abstract}
This document exercises common visual-editor features while preserving ordinary
LaTeX source. Use the editor-title button or run
\texttt{Toggle Visual Editor} from the Command Palette to switch modes.
\end{abstract}

\section{Rich text}
\label{sec:rich-text}

The visual editor should display \textbf{bold text}, \textit{italic text}, \emph{emphasized text}, and a \href{https://www.latex-project.org/}{web link}.
It should also preserve this custom command:

\subsection{Lists and footnotes}

\begin{itemize}
  \item A bulleted item with inline mathematics: \(e^{i\pi} + 1 = 0\).
  \item A second item with a footnote.\footnote{Footnotes are rich LaTeX content.}
\end{itemize}

\begin{enumerate}
  \item First numbered item.
  \item Second numbered item.
\end{enumerate}

\section{Mathematics}

The following displayed equation has a label:
\begin{equation}
  \int_{0}^{1} x^2\,dx = \frac{1}{3}.
  \label{eq:integral}
\end{equation}

Equation~\ref{eq:integral} is referenced without changing its source syntax.

\begin{theorem}
For every real number \(x\), \(x^2 \geq 0\).
\end{theorem}

\section{Table}

\begin{table}[ht]
  \centering
  \begin{tabular}{lc}
    \textbf{Feature} & \textbf{Status} \\
    Headings & Supported \\
    Equations & Supported \\
    Figures & Supported \end{tabular}
  \caption{Content included in this test document.}
  \label{tab:features}
\end{table}


\end{document}
